\section*{Capítulo \ref{ch:b3:milky-way-galaxies} --- A Via Láctea e as galáxias}

\begin{solution}{exo:b3:milky-way-galaxies:1}
(a)~$2\pi R_0 = 2\pi\times\qty{2.523e17}{km} = \qty{1.585e18}{km}$; com
\qty{229}{km.s^{-1}}, isso leva \qty{6.92e15}{s}, ou seja
\qty{219}{Myr}. (b)~$4540/219 = 20.7$: cerca de vinte voltas.
\end{solution}

\begin{solution}{exo:b3:milky-way-galaxies:2}
(a)~$\qty{1}{au}/\qty{1}{yr} = \num{1.496e8}/\num{3.156e7} =
\qty{4.74}{km.s^{-1}}$. (b)~$v = 4.74\times6.411\times8.178 =
\qty{248.5}{km.s^{-1}}$, dentro do intervalo da revisão,
\qty[separate-uncertainty=true]{248 \pm 3}{km.s^{-1}}. (c)~$248.5 - 11 = \qty{237.5}{km.s^{-1}}$,
cerca de \qty{8}{km.s^{-1}} acima dos \qty{229}{km.s^{-1}} das gigantes vermelhas:
os dois dependem de valores distintos de $R_0$ (8.178 e \qty{8.122}{kpc}),
e o próprio \omterm{def:b3:milky-way-galaxies:lsr}{padrão local de repouso} (LSR) pode ter um movimento de até cerca de \qty{15}{km.s^{-1}} em relação
a uma órbita de fato circular.
\end{solution}

\begin{solution}{exo:b3:milky-way-galaxies:3}
(a)~$n = 10(1 - 0.6) = 4$: E4. (b)~$b/a = 1 - 7/10 = 0.3$. (c)~A classificação
usa a razão axial aparente: um elipsoide achatado visto ao longo de seu eixo menor
tem contorno circular e é classificado como E0.
\end{solution}

\begin{solution}{exo:b3:milky-way-galaxies:4}
$d = 10^{\mu/5 + 1}$~pc: a LMC \qty{49.6}{kpc}, Andrômeda
\qty{783}{kpc}, NGC~4258 \qty{7.58}{Mpc}. Uma estrela da Nuvem tem uma
paralaxe de $1/\qty{49.6}{kpc} = \qty{0.020}{mas}$, a ordem de grandeza da melhor
incerteza do catálogo para estrelas brilhantes: uma única paralaxe não fornece
distância alguma, e a Nuvem é situada por outros meios (binárias de
eclipse), a partir das quais começa a escala das cefeidas.
\end{solution}

\begin{solution}{exo:b3:milky-way-galaxies:5}
(a)~$\omega_0 = A - B = \qty{27.2}{km.s^{-1}.kpc^{-1}}$ e $V_0 =
27.2\times8.178 = \qty{222}{km.s^{-1}}$. (b)~$V'_0 = -(A + B) =
\qty{-3.4}{km.s^{-1}.kpc^{-1}}$: uma curva de leve queda. (c)~Ambos estão
abaixo dos valores das gigantes vermelhas (229 e \qty{-1.7}{km.s^{-1}} por
quiloparsec); os autores dessa curva encontraram o mesmo, com $V_0$ cerca de dois
desvios-padrão abaixo, a partir de $R_0$ e dessas \omterm{def:b3:milky-way-galaxies:rotation-curve}{constantes de Oort}. As
fórmulas supõem órbitas circulares em um disco axissimétrico; $C$ e
$K$ não nulos mostram que o campo local de velocidades é perturbado pela barra e pelos
braços espirais, e os valores locais de $A$ e $B$ carregam essas perturbações.
\end{solution}

\begin{solution}{exo:b3:milky-way-galaxies:6}
(a)~$\rho = V^2/(4\pi GR^2) = 229^2/(4\pi\times\num{4.301e-6}\times
8.178^2) = \num{1.45e7}$ massas solares por quiloparsec cúbico, ou
\num{0.0145} por parsec cúbico. (b)~$M_\odot c^2 = \num{1.9884e30}\times
(\num{2.998e8})^2 = \qty{1.787e47}{J} = \qty{1.115e57}{GeV}$, e
$\qty{1}{pc^3} = (\qty{3.0857e18}{cm})^3 = \qty{2.938e55}{cm^3}$: a
razão é \qty{37.96}{GeV.cm^{-3}}. Então $0.30/37.96 = \num{0.0079}$ massas solares
por parsec cúbico. (c)~A densidade da curva plana conta toda a massa
que sustenta as órbitas, estrelas e gás incluídos, como se estivesse espalhada em uma
esfera; a matéria escura é apenas parte dela, cerca de metade aqui.
\end{solution}

\begin{solution}{exo:b3:milky-way-galaxies:7}
(a)~$W_I(\mathrm{LMC}) = -3.314\times1.4771 + 15.888 = \qty{10.993}{mag}$,
logo $M_W = 10.993 - 18.477 = \qty{-7.484}{mag}$. (b)~$\mu = 25.00 + 7.484 =
32.484$ e $d = 10^{32.484/5 + 1}~\mathrm{pc} = \qty{31.4}{Mpc}$.
(c)~$\delta d/d = (\ln 10/5)\,\delta\mu = 0.4605\times0.077 =
\qty{3.5}{\percent}$ para uma estrela, $3.5/\sqrt{50} = \qty{0.5}{\percent}$
para cinquenta, se os erros forem independentes (o erro da distância própria da Nuvem
se soma igualmente a todas elas).
\end{solution}

\begin{solution}{exo:b3:milky-way-galaxies:8}
$\mu = I - A_I - M_I = 26.00 - 0.05 + 4.05 = 30.00$, logo $d =
10^{7}~\mathrm{pc} = \qty{10.0}{Mpc}$. Um erro de \qty{0.04}{mag} equivale a
$0.4605\times0.04 = \qty{1.8}{\percent}$ em distância.
\end{solution}

\begin{solution}{exo:b3:milky-way-galaxies:9}
(a)~$d = cz/H_0$ resulta em $5\log_{10}d = 5\log_{10}(cz) - 5\log_{10}H_0$,
logo $m_B - M_B - 25 = 5\log_{10}(cz) - 5\log_{10}H_0$; com
$\log_{10}(cz) = a_B + 0.2\,m_B$, $5\log_{10}H_0 = 5a_B + M_B + 25$, o que
é o resultado. (b)~$\log_{10}H_0 = -3.8506 + 0.714158 + 5 = 1.8636$, logo
$H_0 = \qty{73.04}{km.s^{-1}.Mpc^{-1}}$. (c)~$M_B = 5(\log_{10}67.4 -
5.714158) = \qty{-19.427}{mag}$: supernovas mais brilhantes em \qty{0.174}{mag},
logo todas as distâncias maiores por um fator $10^{0.2\times0.174} = 1.084$, a
razão $73.04/67.4$. A tensão é uma discordância de oito por cento na
escala de distâncias.
\end{solution}

\begin{solution}{exo:b3:milky-way-galaxies:10}
(a)~Uma estrela na borda, em órbita circular, tem $V^2/R \approx GM/R^2$
(aproximação esférica). (b)~$R = \sqrt{M/(\pi\Sigma)}$, logo $V^2 =
GM/R = G\sqrt{\pi\Sigma}\,M^{1/2}$ e $M = V^4/(\pi\Sigma G^2) \propto
V^4$; com uma razão massa-luminosidade comum, $L \propto V^4$ também.
(c)~$(100/229)^4 = \num{0.036}$: a galáxia mais rápida é 27 vezes mais
massiva. (d)~Mesma luminosidade: o fluxo cai com $1/d^2$, logo a mais fraca
está duas vezes mais distante.
\end{solution}

\begin{solution}{exo:b3:milky-way-galaxies:11}
(a)~$\qty{1}{km.s^{-1}} = \qty{1.0227}{kpc.Gyr^{-1}}$, logo
\qty{229}{km.s^{-1}} é \qty{234.2}{kpc.Gyr^{-1}}. (b)~$2\pi R/V$:
\qty{107}{Myr} a \qty{4}{kpc}, \qty{322}{Myr} a \qty{12}{kpc}.
(c)~$\Delta\phi = Vt\,(1/4 - 1/12) = 234.2\times0.1667 = \qty{39.0}{rad}$,
6.2 voltas: o braço seria uma espiral de seis voltas apertadas entre 4 e
\qty{12}{kpc}, muito mais fechada que os braços observados, e ainda mais a cada
rotação adicional. (d)~Em uma onda de densidade o padrão gira com uma única
velocidade angular, como corpo rígido; as estrelas e o gás atravessam-no, logo a forma
dos braços não se enrola.
\end{solution}

\begin{solution}{exo:b3:milky-way-galaxies:12}
(a)~$a = (4.2\times10^6\times16^2)^{1/3} = \qty{1024}{au}$; a
\qty{8178}{pc}, ele subtende $1024/8178 = \ang{;;0.125}$ (uma au a um parsec
subtende um arco-segundo). (b)~$R_S = 2GM/c^2 = 2\times\num{6.674e-11}
\times\num{8.35e36}/(\num{2.998e8})^2 = \qty{1.24e10}{m} = \qty{0.083}{au}$;
o diâmetro $2R_S$ subtende $0.166/8178 = \num{2.0e-5}$ arco-segundo, cerca de
\qty{20}{\micro as}. (c)~$R_S = \qty{4.02e-7}{pc}$, e $M/(\tfrac43\pi
R_S^3) = \num{1.5e25}$ massas solares por parsec cúbico, quarenta vezes o
limite inferior imposto pelo movimento e pelo tamanho de Sgr~A*: o limite chega
a menos de um fator quarenta da densidade do próprio buraco negro.
\end{solution}

\begin{solution}{pb:b3:milky-way-galaxies:1}

\textbf{Parte I.}

\textbf{1.} $v_r = A\,d\sin2\ell$: $+A d$ em \ang{45}, 0 em \ang{90}, $-A d$ em
\ang{135}, 0 em \ang{180}.

\textbf{2.} $A = (15.4 + 15.2)/(2\times\qty{1.00}{kpc}) =
\qty{15.3}{km.s^{-1}.kpc^{-1}}$.

\textbf{3.} Em \ang{90}, $v_t = d(B - A) = 4.74\times(-5.91)\times1.00 =
\qty{-28.02}{km.s^{-1}}$, logo $B = 15.3 - 28.02 =
\qty{-12.72}{km.s^{-1}.kpc^{-1}}$.

\textbf{4.} Em \ang{0}, $v_t = d(A + B) = 4.74\times0.55 =
\qty{2.61}{km.s^{-1}}$, logo $B = 2.61 - 15.3 = \num{-12.69}$: os dois
concordam dentro de \num{0.03}.

\textbf{5.} $\omega_0 = A - B = \qty{28.0}{km.s^{-1}.kpc^{-1}}$ e
$V_0 = 28.0\times8.18 = \qty{229}{km.s^{-1}}$.

\textbf{6.} $V'_0 = -(A + B) = \qty{-2.6}{km.s^{-1}.kpc^{-1}}$: decrescendo suavemente.

\textbf{7.} $P = 2\pi/\omega_0$: $2\pi\times\qty{8.18}{kpc}/\qty{229.2}{km.s^{-1}}
= \qty{219}{Myr}$; $4540/219 = 20.7$ voltas.

\textbf{Parte II.}

\textbf{8.} O ponto tangente exige uma linha de visada que passe mais perto do
centro do que o Sol está, $\sin\ell > 0$ com $\ell < \ang{90}$ (ou o
quadrante espelhado, $\ang{270} < \ell < \ang{360}$, onde a maior
velocidade é a mais negativa). Fora do círculo solar, a distância
ao centro cresce ao longo de toda a linha de visada: nenhum ponto é
destacado, e a velocidade sozinha não diz onde está o gás.

\textbf{9.} $R = 8.18\sin\ell$: \qtylist{5.26;6.27;7.08;7.69}{kpc}.

\textbf{10.} $V = v_\mathrm{max} + V_0\sin\ell$: $82.9 + 147.2 = 230.1$,
$55.6 + 175.4 = 231.0$, $32.4 + 198.3 = 230.7$, $14.9 + 215.2 =
\qty{230.1}{km.s^{-1}}$.

\textbf{11.} $\Delta\nu = \nu v/c = \num{1.420406e9}\times82.9/299792 =
\qty{393}{kHz}$, em direção a frequências mais baixas: o gás se afasta.

\textbf{12.} Dentro de cerca de \qty{3}{km.s^{-1}} das gigantes vermelhas: entre 5 e
\qty{8}{kpc} a curva é plana.

\textbf{13.} Uma massa pontual dá $V = V_0\sqrt{R_0/R} = 229\times\sqrt{8.18/5.26} =
\qty{286}{km.s^{-1}}$ no primeiro ponto tangente, muito mais rápido do que a
\qty{230}{km.s^{-1}} encontrada: a massa não está concentrada no centro,
grande parte dela está entre 5 e \qty{8}{kpc} e além.

\textbf{Parte III.}

\textbf{14.} $\num{1.067e10} + 2\times\num{3.944e10} = \num{8.96e10}$ massas
solares, quase o dobro do valor da revisão \num[separate-uncertainty=true]{5 \pm 1 e10}; as estrelas do modelo
podem ser pesadas demais, o que faria a matéria escura encontrada abaixo ser uma
subestimativa.

\textbf{15.} $V_b^2 = GMR^2/(R^2 + b^2)^{3/2} = \num{4.301e-6}\times\num{1.067e10}
\times400/(400.09)^{3/2}$, $V_b = \qty{47.9}{km.s^{-1}}$.

\textbf{16.} Disco fino: $(a + b)^2 = 30.80$, $V^2 = \num{4.301e-6}\times\num{3.944e10}
\times400/(430.8)^{3/2}$, $V = \qty{87.1}{km.s^{-1}}$. \omterm{def:b3:milky-way-galaxies:structure}{Disco espesso}:
$(a + b)^2 = 11.56$, $V = \qty{90.1}{km.s^{-1}}$.

\textbf{17.} $V_\mathrm{vis} = \sqrt{47.9^2 + 87.1^2 + 90.1^2} =
\qty{134.2}{km.s^{-1}}$.

\textbf{18.} Em \qty{8.18}{kpc}: 74.8, 108.4 e \qty{127.8}{km.s^{-1}},
no total \qty{183.5}{km.s^{-1}}, frente a $V_0 = \qty{229}{km.s^{-1}}$:
a massa visível já fica aquém no Sol.

\textbf{19.} $\sqrt{G\times\num{8.96e10}/20} = \qty{138.8}{km.s^{-1}}$. O
valor do modelo, \qty{134.2}{km.s^{-1}}, está próximo, porque em \qty{20}{kpc},
quatro a seis comprimentos de escala do disco para fora, quase toda a massa visível está
dentro da órbita: a curva visível tornou-se kepleriana.

\textbf{Parte IV.}

\textbf{20.} $M = 20\times200^2/\num{4.301e-6} = \num{1.860e11}$ massas solares.

\textbf{21.} $M_\mathrm{vis} = 20\times134.2^2/\num{4.301e-6} = \num{8.37e10}$
massas solares.

\textbf{22.} $M_\mathrm{dark} = \num{1.860e11} - \num{8.37e10} = \num{1.02e11}$;
já que ambas as massas são $RV^2/G$ no mesmo $R$, a fração é
$1 - (V_\mathrm{vis}/V)^2 = 1 - (134.2/200)^2 = \num{0.550}$.

\textbf{23.} $x = 20/14.8 = 1.351$, $\ln(2.351) - 1.351/2.351 = 0.8549 - 0.5748
= 0.2801$, $M = 4\pi\times\num{1.06e7}\times14.8^3\times0.2801 =
\num{1.21e11}$ massas solares; a fração escura é
$1.21/(1.21 + 0.837) = \num{0.59}$. O halo é ajustado a toda
a curva, a qual ele superestima a \qty{20}{kpc}: as medições caem ali
cerca de \qty{10}{km.s^{-1}} abaixo do modelo suave (\qty{209}{km.s^{-1}}).

\textbf{24.} $x = 8.18/14.8 = 0.553$, $\rho = \num{1.06e7}/(0.553\times1.553^2)
= \num{7.95e6}$ massas solares por quiloparsec cúbico, \num{0.0080} por
parsec cúbico, ou \qty{0.30}{GeV.cm^{-3}}, o valor que os autores do
ajuste citam.

\textbf{25.} $1 - 8.37/18.60 = \num{0.550}$: a fração de matéria escura da
massa dentro de \qty{20}{kpc} é
\textbf{\qty[text-series-to-math=true]{55}{\percent}}.
\end{solution}
